浅层时空模型已经能够访问本流观测，额外引入局部读出是否值得，仍需匹配对照验证。
本文研究完整历史监督下的稀疏网络流量离线补全，
将一层SPIN时空上下文与本流前后最近真实观测的三统计量读出联合解码。
通过独立训练Context-only和Direct-only，在两套WAN上比较两条路径的使用价值。
固定50槽窗口、20\%输入观测、三种缺失条件和两个配对初始化，
采用120轮训练及预先固定的40轮低学习率诊断，并保留两个预算截面。
160轮预算下，组合相对Context-only的开发NMAE在Abilene和GEANT分别降低1.64\%与2.19\%，
相对Direct-only分别降低2.21\%与4.32\%；两个种子和三个条件方向一致，
末段轨迹也支持这一排序。
相对Context-only，批量8的每窗推理耗时分别增加3.82\%与3.47\%，
但组合明显慢于Direct-only，且并非所有误差指标或观测支持分组都更好。
这些结果支持显式局部读出在当前浅层骨架上的增量，
不证明其引入新信息、优于所有现有方法或具有独立未见流量泛化优势。
本文提供具有明确评价边界的内部方法初稿，公平外部确认仍需完成。
